\scanpage{98}

\[
\unsure{\mathscr S|\mathscr G|\mathfrak S}'_0=\{x\in E\mid w\in W,\ w(\mathbb R x)\ne\mathbb R x\Longrightarrow wx-x\in R\}
\]

\[
\begin{aligned}
x\in\mathscr S'_0,\quad r\in R\qquad&s_rx-x=2\frac{(x,r)}{(r,r)}r\\
s_r(\mathbb R x)&=(\mathbb R x)\Longleftrightarrow s_rx=x\text{ ou }s_rx=-x\\
&\Longleftrightarrow (x,r)=0\text{ ou }x\in\mathbb R r
\end{aligned}
\]

\[(x,r)\ne0\text{ et }x\notin\mathbb R r\Longrightarrow\underbrace{2\frac{(x,r)}{(r,r)}}_{\unsure{n_{x,r}|n_{x r}}}=\pm1\]

\deleted{$(x,r)$ \unsure{de|donc}}\quad .

\diagram{%
\begin{tikzcd}[ampersand replacement=\&,column sep=small,row sep=0pt]
 \mathscr S_0^{\unsure{\prime\prime|\prime}}\&\subset\&\mathscr S'\&
 =\{x\in E\mid \forall r\in R\cap\complement\mathbb R x,\ \text{on a }(x,r)\in\tfrac{(r,r)}2\{-1,0,+1\}\}\\
 \cup\&\&\cup\\
 \mathscr S_0\&\subset\&\mathscr S
\end{tikzcd}%
}

\[
\begin{aligned}
\mathscr S'&=\mathscr S\cup\{\lambda s\mid s\in R,\lambda\in\mathbb R^{++},\
 (\lambda s,\unsure{r})\in\tfrac{(r,r)}2\{-1,0,1\}\\
&\hspace{16em}\forall r\in R-\{s,-s\}\}\\
&\text{i.e.}\qquad \forall r\in R-\{s,-s\}\unsure{-|\ne}(R\cap s^\perp)\\
&\deleted{\text{i.e. }\lambda\in\tfrac{(r,r)}{2(s,r)}}\qquad
 \lambda\in\tfrac{(r,r)}{2(s,r)}\{-1,0,1\}\\
&\text{i.e. }\lambda=\pm\tfrac{(r,r)}{2(s,r)}\qquad
 \text{i.e. }\tfrac1\lambda=\pm\tfrac{2(s,r)}{(r,r)}=\pm n(s,r)
\end{aligned}
\]

un tel $\lambda$ existe (il est unique au signe près), pour $s$ donné, ssi les quantités $\left|\frac{(s,r)}{(r,r)}\right|$, pour $r$ racine non orthogonale ni prop. à $s$, sont égales, i.e. ssi il n’y a qu’un seul angle de droites possible (savoir $\frac\pi3$, \inserted{d’où $n(s,r)=\pm1$, d’où $\lambda=\unsure{s|1|\pm1}$}) pour des racines distinctes non orth. C’est le cas ssi ttes les racines sont de $=$ longueur, i.e. les types $A_n$, $D_n$, $E_i$ $(6\leq i\leq8)$.

% archive sheet number: 99
